
Repository maintenance prediction has practical importance for software supply chain security and dependency management. Prior work has employed graph-based centrality analysis requiring distributed infrastructure \citep{he2024repository} and manually-tuned composite metrics \citep{adejumo2025stability} without establishing whether simpler methods achieve comparable performance. This study tests whether logistic regression trained on basic GitHub metadata features can achieve $\geq 75$\% accuracy for binary maintenance classification, addressing a gap in the literature where complex methods were deployed without simple baseline comparisons.

He et al. (2024) achieved C-Index 0.810 (approximately 80-85\% classification accuracy) on 103,354 GitHub repositories using gradient boosting with HITS centrality features, requiring expensive graph construction. Adejumo and Johnson (2025) proposed a Composite Stability Index (CSI) aggregating repository metrics with manually-tuned weights (30\% activity, 25\% commits, 25\% issues, 20\% age), achieving F1 0.80 on 100 repositories. Neither study tested whether logistic regression on basic metadata could match or exceed these results.

The present study evaluated logistic regression on 120 real Papers with Code benchmark repositories using six core GitHub API features: stars, forks, contributors, commits, open issues, and days since last commit. Binary labels were assigned using a 180-day threshold (maintained if days\_since\_last\_commit < 180, abandoned otherwise). Initial experiments with eight features (including two derived features: commit frequency and issue resolution rate) achieved 100\% test accuracy but were invalidated due to tautological relationships in the derived features. Subsequent experiments restricted to six independently measured GitHub features achieved 95.8\% logistic regression accuracy and 100\% gradient boosting accuracy on a 24-sample held-out test set.

Contributions of this work are as follows. First, it establishes an empirical simplicity baseline: logistic regression achieves 95.8--100\% accuracy on Papers with Code benchmark repositories, exceeding the 75\% target threshold. Second, it reveals a two-tier signal hierarchy where staleness (days\_since\_last\_commit) provides the dominant signal (coefficient -3.05) relative to engagement features (+0.14 to +0.55). Third, it quantifies the marginal value of ensemble complexity: gradient boosting provides 4.2 percentage points improvement over logistic regression. Fourth, it demonstrates that six GitHub API features suffice without graph-based centrality features.

Scope limitations include domain restriction (Papers with Code machine learning benchmark repositories only), small sample size (120 repositories, wide confidence intervals), and absence of temporal validation (stratified IID split only, no temporal holdout). Initial experimental validity issues (tautological derived features) were detected and corrected, with final results reported on six real GitHub metadata features.

The remainder of this paper is organized as follows. Section 2 reviews related work on repository maintenance prediction and positions the current contribution. Section 3 describes the methodology including dataset collection, feature engineering, and model training. Section 4 details the experimental protocol and validation gates. Section 5 presents classification performance, coefficient analysis, and complexity comparison results. Section 6 discusses interpretation, limitations, and positioning relative to prior work. Section 7 concludes.
